\documentclass[11pt]{article}
 
% --- Page and fonts ---
\usepackage[margin=1in]{geometry}
\usepackage[T1]{fontenc}
\usepackage[utf8]{inputenc}
\usepackage{lmodern}
\usepackage{microtype}
 
% --- Citations (author-year, as in last year's proposal) ---
\usepackage[authoryear,round]{natbib}
 
% --- Misc ---
\usepackage{hyperref}
\hypersetup{colorlinks=true, linkcolor=black, citecolor=black, urlcolor=black}

\title{Knuckleball: Auditable Entropy Provenance and Exact Replay for Stochastic Scientific Computation\\[0.4em]\large ICSRI Research Proposal}
\author{Nikhil Sunder\thanks{ICSRI Research Fellow, Intelligent Computer Systems Research Institute, Miami Herbert Business School, University of Miami. Email: \texttt{nss106@miami.edu}. ORCID: 0009-0007-3323-1760.}\\\normalsize Miami Herbert Business School, University of Miami}
\date{}

\begin{document}
\maketitle

\begin{abstract}
Stochastic scientific computation in Python has no provenance layer: randomness enters
through an opaque system call, is expanded by a pseudo-random generator, and is
reproduced only by re-supplying a seed, which records nothing about the entropy source,
cannot detect its degradation, and does not guarantee bit-exact replay across platforms.
This proposal outlines the research proposition for Knuckleball: a Python entropy layer
for stochastic scientific computation that harvests randomness from physical and system
sources, conditions it under an explicit min-entropy budget, monitors source health with
standardized statistical tests, and records emitted randomness in a verifiable
transcript so that non-deterministic computations can be replayed exactly and their
randomness audited by third parties.
\end{abstract}

\noindent\textbf{Broad Topic:} Systems Design and Systems Research
 
\noindent\textbf{Narrowed Topic:} Entropy provenance, source health, and exact replay for pseudo-random and physical randomness in scientific computation
 
\noindent\textbf{Final Artifacts:} Python package (PyPI, conda-forge, OpenSSF Best Practices); research paper and JOSS software paper; reproducible notebooks and cross-platform test harness

\section{Introduction}
 
Every random number produced by a computer passes through a three-stage pipeline
\citep{turan2018,barker2015}. A noise source exposes an unpredictable physical
process as raw bits that are biased and correlated; a conditioner compresses them into a
shorter, near-uniform block \citep{vonneumann1951,impagliazzo1989}; a deterministic
random bit generator (DRBG) expands that block into an arbitrarily long stream
\citep{matsumoto1998,oneill2014}. The quantity flowing through the pipeline is
min-entropy, the worst-case unpredictability of the source.
 
In scientific computation the pipeline is invisible. The operating system runs all three
stages inside one system call, and the practitioner either consumes its output directly
or bypasses the first two stages by supplying a fixed seed. Bayesian estimation
\citep{geweke1999}, bootstrap inference \citep{kilian1998}, simulation-based testing
\citep{diebold1995}, and reinforcement learning \citep{henderson2018} all obtain
randomness this way, and the seed becomes the entire record of where it came from.
 
A seed reproduces the expansion stage and nothing else. It records nothing about which
sources contributed or whether they were functioning, it cannot detect a generator that
has silently begun returning constants, and it does not guarantee bit-exact
reproduction, because the expanded stream depends on the DRBG implementation as well as
the seed. Seed choice alone produces reported-effect-sized variation in machine
learning \citep{henderson2018,agarwal2021}, and computational reproducibility remains a
first-order concern across the quantitative sciences \citep{peng2011,stodden2016}. No
Python tool addresses the gap: \texttt{os.urandom} is opaque by design, NumPy's
\texttt{SeedSequence} handles expansion only \citep{harris2020}, and randomness beacons
provide public values without local sources or health monitoring
\citep{syta2017,kelsey2019}.
 
This proposal asks whether entropy provenance, source integrity, and exact replay can be
made first-class properties of stochastic scientific computation, and what system design
achieves them. The proposed system is Knuckleball, a Python entropy layer intended to
credit min-entropy from heterogeneous sources under explicit assertions, condition it
within the credited budget, monitor source health with standardized tests, and record
emitted randomness in a tamper-evident transcript for replay, while leaving existing
consumers of \texttt{random.Random} and NumPy unchanged.
 

\section{Problem and Design Goals}
 
The provenance gap decomposes into three problems. \emph{Provenance}: for any block of
randomness, which sources contributed, how much entropy each was credited, and when.
\emph{Integrity}: whether a source is still delivering the entropy it is credited with,
given that a failed source returns bytes of the correct type and raises no error.
\emph{Replay}: whether a third party can reconstruct the randomness a computation
consumed on different hardware and library versions, without trusting the original
machine.
 
Four failure modes with documented precedent motivate the design \citep{turan2018}:
silent source death, where a generator returns constants after a fault; over-extraction,
where a conditioner emits more bits than its input entropy justifies; over-trusted
assertion, where a source is credited at a rate it does not deliver; and implementation
drift, where a fixed seed yields different streams across library versions. An active
adversary who observes or influences the source at run time is out of scope; that is the
threat model of cryptographic generation, which the system does not address.
 
The design is organized around four goals, stated here as targets rather than as
guarantees, since their exact form will be settled by implementation.
 
\begin{itemize}
  \item \textbf{G1 (Budget accounting).} Conditioned output should be bounded by the
        min-entropy credited from sources, following the extraction bounds established
        in the theoretical literature, so that the system never silently degrades into a
        pseudo-random generator while presenting itself as physical.
  \item \textbf{G2 (Health monitoring).} Sources should be monitored continuously by
        standardized health tests, excluded from crediting when they fail, and auditable
        against their asserted rates by empirical estimation.
  \item \textbf{G3 (Replay).} Emitted randomness should be recorded with its provenance
        in a tamper-evident form, such that a computation can be re-run on the recorded
        randomness rather than on live sources.
  \item \textbf{G4 (Transparency).} Existing consumers of Python's and NumPy's random
        interfaces should operate unchanged, with the provenance layer confined to the
        seeding step and adding no statistical distortion.
\end{itemize}
 
Three exclusions are explicit: the system is not a cryptographic generator and must not
be used for keys or tokens; it does not stream physical entropy into every variate but
draws it to seed an expander; and the initial system is pure Python, with hardware
instructions and sensor sources as optional extras.
 

\section{Architecture}
 
The anticipated architecture follows the pipeline of Section~1 directly, with one
layer per stage and a recording layer alongside. A \emph{source} layer wraps a physical
or system process and exposes raw bytes together with a declared entropy rate, a flag
indicating whether the output is public, and a flag indicating whether acquisition
performs I/O. The declared rate is the one place a human judgement enters the pipeline,
and the design records it rather than inferring it. A \emph{health} layer applies the
SP~800-90B repetition-count and adaptive-proportion tests to each source's raw output
and marks sources that fail. A \emph{conditioning} layer compresses raw bytes to
near-uniform output; a universal-hash extractor is the expected default, with classical
debiasing and the SP~800-90A derivation function as alternatives. A \emph{pool} owns
the entropy ledger, coordinating credit from healthy sources, conditioning, debit, and
emission, with a configurable policy for reads the ledger cannot cover. A
\emph{bridge} layer adapts the pool to \texttt{random.Random} and to NumPy's
\texttt{SeedSequence}, and a \emph{transcript} records emitted randomness with a hash
chain and stands in for the sources during replay.
 
The intended consumption pattern is seed-then-expand: one conditioned block per
computation seeds an existing expander, and the transcript holds one record per run.
This confines overhead to a single read and keeps the provenance layer invisible past
the seed, which is what makes G4 achievable without modifying downstream code. Streaming
physical entropy into every variate would offer no statistical benefit to simulation and
would make replay impractical.
 
Two conventions carry over from the author's existing packages. Asynchronous interfaces
will exist only where I/O occurs, principally for network beacon sources. And no caching
may sit on the entropy path, since a cached block re-emitted represents zero new entropy
credited as fresh; this will be enforced by test rather than by convention.
 
The public interface, module boundaries, and configuration surface are deliberately
left open at this stage. Decisions on which conditioners ship by default, how the
ledger handles concurrent readers, and how transcripts are serialized will be made
against working code and the workloads of Section~6.

\section{Theoretical Basis}
 
The design rests on established results rather than new theory; this section identifies
them and the role each plays.
 
\textbf{Min-entropy and extraction.} Min-entropy measures the probability of the single
most likely outcome and is the correct quantity for randomness extraction, because it
bounds the success of the best guess rather than the average surprise. The leftover
hash lemma \citep{impagliazzo1989,hastad1999} establishes that a universal hash family
applied to a source of known min-entropy yields output close to uniform provided the
output is shorter than the input entropy by a margin that depends on the desired
closeness. This result sets the accounting rule for G1: the ledger debits output length
plus the lemma's loss term. Toeplitz matrices are a standard universal family with a
short seed and an elementary universality argument \citep{krawczyk1994}, which is why
they are the expected default conditioner.
 
\textbf{Entropy estimation and health testing.} NIST SP~800-90B \citep{turan2018}
specifies both a battery of estimators for min-entropy on non-IID sources and two
continuous health tests: a repetition-count test that detects a source stuck on a
constant, and an adaptive-proportion test that detects a source whose most likely value
has become too frequent. Both tests are parameterized by the source's asserted rate and
a target false-alarm probability, and both have bounded detection latency under those
parameters. G2 adopts the standard's tests and estimators directly rather than
designing new ones.
 
\textbf{Hash chaining.} Linking each transcript record to its predecessor through a
collision-resistant hash makes any modification, deletion, or truncation detectable on
verification, a construction with long precedent in timestamping and audit logging.
G3 uses it to make transcripts tamper-evident; the same construction underlies the
NIST beacon format \citep{kelsey2019}.
 
\textbf{Data processing.} Statistical distance cannot increase under a deterministic
function. Because an expander is deterministic in its seed, a consumer fed a seed that
is close to uniform sees output that is at least as close to what it would see from an
ideal seed. This is the basis for G4: transparency follows from the seed's quality and
requires nothing of the consumer.

\section{Positioning}
 
Knuckleball composes standardized components rather than replacing them. NIST
SP~800-90A, B, and C \citep{barker2015,turan2018,barker2024} together define the
source, conditioner, and expander stages and the rules for composing them; the proposed
system implements the source and conditioner stages in Python with the standard's own
health tests and estimators, and delegates expansion to existing DRBGs.
 
Existing Python tools each cover one stage. \texttt{os.urandom} and \texttt{secrets}
expose the operating system's full pipeline as an opaque call, which is correct for
cryptography and uninformative for provenance. NumPy's \texttt{SeedSequence} and
\texttt{Generator} \citep{harris2020} provide high-quality expansion and seed spreading
but begin at the seed. Randomness beacons such as drand and the NIST Beacon
\citep{syta2017,kelsey2019} publish verifiable public values on a schedule, which
serves public randomization but offers no local source, no conditioning, and no health
monitoring. Hardware generator interfaces expose a single vendor source at an asserted
rate with no accounting. To the author's knowledge no tool in the Python ecosystem
makes the source, the entropy budget, the health state, and the emitted stream jointly
visible, and none provides replay independent of the expander implementation.
 
Python and scientific computation are the chosen scope because that is where the
reproducibility burden falls in practice: the workloads of Section~6 are written in
Python against NumPy, and their results are expected to be reproducible by reviewers
who do not share the author's hardware.

\section{Application}
 
The system will be attached to three standard econometric workloads chosen to span
Bayesian estimation, resampling inference, and simulation-based testing. In each case
the attachment is the same: one conditioned block seeds the workload's generator, the
transcript records that block with its provenance and health state, and a second
machine replays the transcript and re-runs the workload.
 
\begin{itemize}
  \item \textbf{Bayesian vector autoregression.} A BVAR with a Minnesota prior
        \citep{litterman1986} estimated by Gibbs sampling on a standard macroeconomic
        dataset, where posterior forecast bands depend on the sampler's random draws.
  \item \textbf{Bootstrap impulse-response intervals.} Bias-corrected bootstrap
        confidence intervals for impulse responses \citep{kilian1998}, where interval
        endpoints depend on the resampling sequence.
  \item \textbf{Monte Carlo evaluation of a forecast comparison test.} The empirical
        size of the Diebold--Mariano test \citep{diebold1995} under a known data-generating
        process, where rejection rates depend on the simulated error sequences.
\end{itemize}
 
Each workload demonstrates three things: that the provenance layer attaches without
changes to the estimation code (G4), that the run replays exactly on a second platform
(G3), and, as an illustration rather than a finding, how wide the band of results is
across independently seeded runs.

\section{Evaluation}
 
\subsection{Validation of design goals}
 
Each design goal will be validated by a procedure that could fail.
 
\begin{itemize}
  \item \textbf{G1.} Conditioned output will be subjected to standard statistical test
        batteries \citep{lecuyer2007}, and the ledger will be audited to confirm that
        no read exceeded the credited budget under the accounting rule of Section~4.
  \item \textbf{G2.} Failures will be injected into each source---constant output,
        stuck bits, collapsed variance---and detection latency and false-alarm rate
        will be measured against the bounds implied by the SP~800-90B parameters.
        Asserted rates for each source will be compared against estimated rates.
  \item \textbf{G3.} Transcripts recorded on macOS (ARM) will be replayed on Linux (x86)
        and Windows, and the emitted streams compared byte for byte. Modified and
        truncated transcripts will be verified to fail.
  \item \textbf{G4.} Each workload of Section~6 will be run with and without the
        provenance layer; output distributions will be compared and wall-clock overhead
        measured.
\end{itemize}
 
\subsection{Metrics}
 
Test-battery pass rates on conditioned output; budget-audit violations (target zero);
detection latency in samples and false alarms per healthy sample; exact-match rate
across platforms; overhead as a fraction of workload wall-clock time; and, for each
workload, the dispersion of results across independently seeded runs.
 
\subsection{Hypotheses}
 
\begin{enumerate}
  \item \textbf{Accounting hypothesis:} budgeted output passes the same statistical
        batteries as the underlying expander while never exceeding credited entropy.
  \item \textbf{Detection hypothesis:} every injected failure mode is detected within
        the latency implied by the health-test parameters, with no false alarms over an
        extended healthy run.
  \item \textbf{Replay hypothesis:} every workload replays exactly across the three
        platforms.
  \item \textbf{Cost hypothesis:} overhead is confined to seeding and is a small fraction
        of workload wall-clock time.
\end{enumerate}
 
\subsection{Deliverables}
 
A Python package published to PyPI and conda-forge under OpenSSF Best Practices; a
written specification of the design and its accounting rules; the fault-injection and
cross-platform test harness; reproducible notebooks for the three workloads; a JOSS
software paper; and a full-length paper describing the design and its validation for a
computational-methods or reproducibility venue.

\bibliographystyle{plainnat}
\bibliography{references}
 
\end{document}
